\begin{figure}[H]
\centering
\begin{tikzpicture}[>=Latex,scale=0.96]
    % Estado inicial
    \node[anchor=west] at (-5.35,4.72)
        {\textbf{(a) Evento inicial: $t=t'=0$}};

    \draw[->,black,very thick] (-4.8,3.62) -- (4.9,3.62)
        node[right] {$x$};
    \draw[->,blue!75!black,thick] (-4.8,3.36) -- (4.9,3.36)
        node[right] {$x'$};

    \draw[black,very thick] (0,3.48) -- (0,4.0);
    \node[above=2pt] at (0,4.0) {$O=O'$};

    \draw[red!75!black,very thick,fill=red!12] (0,3.49) circle (0.17);
    \foreach \ang in {0,45,...,315}
        \draw[red!75!black,thick]
            (0,3.49) ++(\ang:0.23) -- ++(\ang:0.16);
    \node[gray!70!black,font=\scriptsize] at (-2.15,4.05)
        {orígenes coincidentes};
    \draw[->,gray!70!black,thin] (-1.35,3.96) -- (-0.25,3.62);
    \node[red!75!black,font=\small,align=center] at (2.35,4.02)
        {emisión del pulso\\de luz};
    \draw[->,gray!70!black,thin] (1.75,3.86) -- (0.28,3.58);

    % Estado posterior
    \node[anchor=west] at (-5.35,2.88)
        {\textbf{(b) Instante posterior}};

    \coordinate (O) at (-4.55,1.52);
    \coordinate (Op) at (-0.75,1.52);
    \coordinate (P) at (4.65,1.52);

    % Ejes espaciales de cada marco
    \draw[->,black,very thick] (-5.05,1.52) -- (5.15,1.52)
        node[right] {$x$};
    \draw[->,blue!75!black,thick] (-1.2,1.25) -- (5.15,1.25)
        node[right] {$x'$};

    \draw[black,very thick] (O) -- ++(0,0.52);
    \node[above=2pt] at ($(O)+(0,0.52)$) {$S$};
    \node[below=5pt] at (O) {$O$};

    \draw[blue!75!black,very thick] ($(Op)+(0,-0.27)$) -- ++(0,0.79);
    \node[blue!75!black,above=2pt] at ($(Op)+(0,0.52)$) {$S'$};
    \node[blue!75!black,below=5pt] at ($(Op)+(0,-0.27)$) {$O'$};
    \draw[->,blue!75!black,very thick]
        ($(Op)+(0.35,0.78)$) -- ++(2.1,0)
        node[midway,above] {$\vec v$};

    % Pulso en el evento P
    \draw[red!75!black,very thick,fill=red!12] (P) circle (0.17);
    \foreach \ang in {0,45,...,315}
        \draw[red!75!black,thick]
            (P) ++(\ang:0.23) -- ++(\ang:0.16);
    \node[red!75!black,font=\small,align=center] at (4.25,2.17)
        {evento $P$};

    % Guías verticales
    \draw[gray!60,densely dotted] (O) -- ($(O)+(0,-2.15)$);
    \draw[gray!60,densely dotted] (Op) -- ($(Op)+(0,-2.15)$);
    \draw[gray!60,densely dotted] (P) -- ($(P)+(0,-2.15)$);

    % Mediciones espaciales
    \draw[<->,black,thick]
        ($(O)+(0,-0.82)$) -- ($(P)+(0,-0.82)$)
        node[midway,above=3pt] {$x=ct$};
    \draw[<->,blue!75!black,thick]
        ($(Op)+(0,-1.38)$) -- ($(P)+(0,-1.38)$)
        node[midway,above=3pt] {$x'=ct'$};
    \draw[<->,gray!75!black,thick]
        ($(O)+(0,-1.94)$) -- ($(Op)+(0,-1.94)$)
        node[midway,above=3pt] {$vt$};

    \node[gray!70!black,font=\scriptsize] at (0,-0.88)
        {esquema espacial; las escalas de ambos marcos se indican de forma conceptual};
\end{tikzpicture}
\caption{Observación de un mismo pulso de luz desde dos marcos inerciales. En el evento inicial coinciden los orígenes y se emite el pulso. Posteriormente, $O'$ se ha desplazado respecto a $O$, mientras ambos observadores asignan al evento $P$ las posiciones $x=ct$ y $x'=ct'$.}
\label{fig:calibrador-luz}
\end{figure}